\chapter{Geometryzacja rozmaitości trójwymiarowych}
\label{ch:geometrizacja-thurstona}
\index{geometryzacja Thurstona}
\index{geometrie Thurstona}

Rozmaitość trójwymiarowa może być złożona z obszarów o odmiennej
geometrii. Torus $T^3$ ma metrykę płaską, $S^3$ metrykę o dodatniej
krzywiźnie stałej, a istnieją zamknięte rozmaitości o krzywiźnie
stałej ujemnej. Twierdzenie geometryzacyjne mówi, w jakim sensie
takie modele wystarczają po rozcięciu dowolnej zamkniętej
trójrozmaitości. Poniżej podamy jego dokładną wersję topologiczną,
opiszemy osiem geometrii i wskażemy miejsce potoku Ricciego w dowodzie.
Głębokie wyniki przywołujemy ze źródeł; nie dowodzimy tutaj
twierdzenia geometryzacyjnego.

Przez \emph{zamkniętą} rozmaitość rozumiemy zwartą rozmaitość bez
brzegu. W całym rozdziale rozmaitości są gładkie i spójne, a w
sformułowaniu geometryzacji również orientowalne. Geometrie Kleina,
zdefiniowane przez jednorodne przestrzenie $G/H$, omówiliśmy szerzej
w Rozdziale~\ref{ch:geometria-kleina}. Tu potrzebny jest jeszcze
warunek metryczny: grupa symetrii ma zachowywać pewną metrykę
riemannowską.

\section{Od geometrii lokalnej do ilorazu}
\label{sec:thurston-model}

\begin{definicja}[Geometria modelowa]
\label{def:thurston-model}
\index{geometria modelowa}
\emph{Geometrią modelową wymiaru trzy} nazywamy parę $(X,G)$,
w której $X$ jest jednospójną gładką trójrozmaitością, a grupa Liego
$G$ działa na $X$ gładko, efektywnie i przechodnio, ze zwartym
stabilizatorem punktu. Wymagamy ponadto, by istniała
dyskretna podgrupa $\Gamma<G$, działająca swobodnie i taka, że
$X/\Gamma$ jest zwarty. Geometria jest \emph{maksymalna}, gdy
działania $G$ nie da się właściwie powiększyć do działania o tych
samych własnościach na $X$.

Zwartość stabilizatora umożliwia wybór $G$-niezmienniczej metryki
riemannowskiej: wybieramy iloczyn skalarny w jednej przestrzeni
stycznej, uśredniamy go po stabilizatorze i przenosimy działaniem
$G$ do pozostałych punktów. Przykładem jest
$(\R^3,\R^3\rtimes\mathrm O(3))$ z metryką euklidesową.
\end{definicja}

W tym rozdziale \emph{rozmaitość geometryczna} oznacza zamknięty
iloraz $M=X/\Gamma$ maksymalnej geometrii modelowej; $\Gamma$ jest
dyskretna, działa swobodnie i zachowuje metrykę. Metryka na $M$
jest lokalnie taka sama jak na $X$, a $X$ jest nakryciem uniwersalnym
$M$. Dla niezwartej części rozmaitości po cięciu użyjemy osobno
sformułowania o zupełnej metryce hiperbolicznej skończonej objętości.
Każdy powyższy model ma metrykę zupełną: jednorodna metryka
riemannowska jest zupełna.

Warunek maksymalności ma znaczenie. Na $\R^3$ same translacje także
działają przechodnio, lecz ich działanie rozszerza się do pełnej
grupy izometrii euklidesowych. Osiem geometrii Thurstona to osiem
\emph{maksymalnych} geometrii modelowych w wymiarze trzy, które mają
zwarty iloraz. Klasyfikacja tej listy jest wynikiem zewnętrznym;
zob. \href{https://doi.org/10.1112/blms/15.5.401}
{P.\ Scott, \emph{The Geometries of 3-Manifolds},
Bull. London Math. Soc. 15 (1983), 401--487, \S\,4}.

Para $(X,G)$ opisuje możliwe lokalne symetrie. Torus $T^3$ jest
\emph{ilorazem} modelu euklidesowego, a nie dziewiątym modelem.
Podobnie $\mathbb{RP}^3$ jest ilorazem $S^3/\{\pm1\}$ i lokalnie
ma geometrię sferyczną, choć może też wystąpić jako odmienna
przestrzeń jednorodna w geometrii rzutowej Kleina.

\section{Osiem geometrii}
\label{sec:thurston-eight}

Podajemy standardową metrykę lub wygodny jej przedstawiciel.
Zmiana wspólnej skali nie zmienia typu geometrii. Przez
$\operatorname{Isom}(X)$ oznaczamy grupę izometrii, a przez
$\operatorname{Isom}_0(X)$ jej składową jedności. Znaki krzywizn
poniżej dotyczą krzywizny sekcyjnej; w trzech ostatnich modelach
zależy ona od płaszczyzny stycznej. Opis i klasyfikacja są zgodne
z cytowaną pracą Scotta oraz z
\href{https://library.slmath.org/nonmsri/gt3m/}
{W.\ P.\ Thurstonem, \emph{The Geometry and Topology of
Three-Manifolds}, zwłaszcza rozdz.\ 1--3}.

\subsection*{Trzy geometrie stałej krzywizny}

\textbf{Geometria sferyczna $S^3$.}
\index{geometrie Thurstona!sferyczna}
Sfera jednostkowa w $\R^4$ ma $K=1$ i pełną grupę izometrii
$\mathrm O(4)$, wymiaru $6$. Zamknięte ilorazy $S^3/\Gamma$,
gdzie skończona grupa $\Gamma<\mathrm O(4)$ działa swobodnie,
nazywają się \emph{sferycznymi formami przestrzennymi}.
Ich grupa podstawowa jest skończona.
Konkretnymi przykładami są przestrzenie soczewkowe
$L(p,q)$: dla $p\geq2$, $\gcd(p,q)=1$, grupa cykliczna rzędu $p$
działa na $S^3\subset\C^2$ przez
\[
 (z_1,z_2)\longmapsto
 (e^{2\pi i/p}z_1,e^{2\pi iq/p}z_2).
\]
Żadna nietrywialna potęga tego przekształcenia nie ma punktu
stałego: co najmniej jedna współrzędna jest niezerowa, a
$\gcd(p,q)=1$. Dla $p=2,q=1$ dostajemy $\mathbb{RP}^3$.

\textbf{Geometria euklidesowa $\mathbb E^3$.}
\index{geometrie Thurstona!euklidesowa}
Jest to $\R^3$ z $K=0$ i
$\operatorname{Isom}(\mathbb E^3)=\R^3\rtimes\mathrm O(3)$,
również wymiaru $6$. Dla kraty $\Lambda\cong\Z^3$ iloraz
$\R^3/\Lambda=T^3$ jest płaskim torusem. Twierdzenie Bieberbacha,
przywołane tu jako wynik zewnętrzny, mówi, że grupa podstawowa
każdej zamkniętej płaskiej trójrozmaitości zawiera $\Z^3$ jako
podgrupę skończonego indeksu. Źródło: cytowany Scott, \S\,4.

\textbf{Geometria hiperboliczna $\mathbb H^3$.}
\index{geometrie Thurstona!hiperboliczna}
W modelu górnej półprzestrzeni $z>0$ można wziąć metrykę
\[
 g=\frac{(\dd x)^2+(\dd y)^2+(\dd z)^2}{z^2},\qquad K=-1.
\]
Pełna grupa izometrii jest izomorficzna z grupą przekształceń
Lorentza zachowujących dodatnią powłokę hiperboloidy
w $\R^{3,1}$; jej wymiar to $6$. Istnieją zarówno zamknięte
ilorazy $\mathbb H^3/\Gamma$, jak i niezwarte ilorazy zupełne
o skończonej objętości. Do drugich należy dopełnienie węzła
ósemkowego w $S^3$; jego końce są \emph{ostrzykami}, których
przekroje są torusami. Istnienie takiej metryki i twierdzenie
o hiperbolicznych wypełnieniach Dehna pochodzą od Thurstona
(\href{https://library.slmath.org/nonmsri/gt3m/}
{\emph{The Geometry and Topology of Three-Manifolds}, rozdz.\ 4}).
Grupa podstawowa zamkniętego ilorazu jest nieskończona i nie jest
wirtualnie abelowa, w odróżnieniu od przypadku euklidesowego.

\subsection*{Dwa iloczyny powierzchni z prostą}

\textbf{Geometria $S^2\times\R$.}
\index{geometrie Thurstona!$S^2\times\R$}
Metryka iloczynowa ma krzywiznę $1$ na płaszczyznach stycznych
do $S^2$ i $0$ na płaszczyznach zawierających kierunek $\R$.
Grupą izometrii jest
$\operatorname{Isom}(S^2)\times\operatorname{Isom}(\R)$,
o składowej jedności wymiaru $4$.
Iloraz $S^2\times(\R/\Z)=S^2\times S^1$ jest zamknięty i ma
grupę podstawową $\Z$. Ten przykład jest \emph{pierwszy}, lecz
nie \emph{nierozkładalny}: sfera $S^2\times\{t\}$ nie ogranicza
kuli. Innym orientowalnym ilorazem tej geometrii jest
$\mathbb{RP}^3\#\mathbb{RP}^3$. Ta rozmaitość nie jest pierwsza:
jej dwa składniki pierwsze mają geometrię sferyczną.

\textbf{Geometria $\mathbb H^2\times\R$.}
\index{geometrie Thurstona!$\mathbb H^2\times\R$}
Metryka iloczynowa ma krzywiznę $-1$ na płaszczyznach stycznych
do $\mathbb H^2$ i $0$ na mieszanych. Grupa izometrii jest
$\operatorname{Isom}(\mathbb H^2)\times\operatorname{Isom}(\R)$,
ze składową jedności wymiaru $4$.
Jeżeli $\Sigma_g$ jest zamkniętą orientowalną powierzchnią
rodzaju $g\geq2$ z metryką hiperboliczną, to
$\Sigma_g\times S^1$ jest zamkniętym ilorazem tej geometrii.
Jego grupa podstawowa to $\pi_1(\Sigma_g)\times\Z$;
drugi czynnik opisuje regularne włókno okręgowe.
Istnienie metryki na każdej zamkniętej orientowalnej
$\Sigma_g$ o $g\geq2$ jest częścią twierdzenia o uniformizacji,
tu użytego jako wynik zewnętrzny; zob.
\href{https://library.slmath.org/nonmsri/gt3m/}
{W.\ P.\ Thurston, \emph{The Geometry and Topology of
Three-Manifolds}, rozdz.\ 2}.

\subsection*{Trzy geometrie związane z grupami Liego}

\textbf{Geometria $\widetilde{\mathrm{SL}}(2,\R)$.}
\index{geometrie Thurstona!$\widetilde{\mathrm{SL}}(2,\R)$}
Przestrzeń stycznych wektorów jednostkowych do
$\mathbb H^2$ jest izomorficzna z $\mathrm{PSL}(2,\R)$;
jej nakrycie uniwersalne daje model
$\widetilde{\mathrm{SL}}(2,\R)$.
Podnosimy nań metrykę Sasakiego. Naturalne odwzorowanie
na $\mathbb H^2$ ma jednowymiarowe włókna, a składowa
jedności grupy izometrii ma wymiar $4$.
Krzywizna sekcyjna przyjmuje znaki dodatnie i ujemne;
nie jest to przestrzeń stałej krzywizny.
Dla hiperbolicznej powierzchni $\Sigma_g$ jej wiązka
jednostkowych wektorów stycznych $U\Sigma_g$ jest zamkniętym
ilorazem tego modelu. Jako wiązka okręgowa nad $\Sigma_g$
ma niezerową liczbę Eulera $\chi(\Sigma_g)=2-2g$.
To właśnie różni ją od iloczynu $\Sigma_g\times S^1$.

\textbf{Geometria $\mathrm{Nil}$.}
\index{geometrie Thurstona!Nil}
Wygodny model to rzeczywista grupa Heisenberga
\[
 N=\left\{\begin{pmatrix}1&x&z\\0&1&y\\0&0&1\end{pmatrix}
 :x,y,z\in\R\right\},
 \qquad
 g=(\dd x)^2+(\dd y)^2+(\dd z-x\,\dd y)^2.
\]
Metryka jest niezmiennicza względem lewych translacji.
W ortonormalnej bazie lewostronnie niezmienniczej $X,Y,Z$
jedyny niezerowy nawias bazowy to $[X,Y]=Z$. Wzór na koneksję
metryczną daje $K(X,Y)=-3/4$ oraz
$K(X,Z)=K(Y,Z)=1/4$. Pełna składowa jedności grupy
izometrii ma wymiar $4$: do translacji dochodzą obroty
płaszczyzny rozpiętej przez $X,Y$.
Macierze o całkowitych $x,y,z$ tworzą kratę $N_{\Z}$,
bo przy mnożeniu nowa współrzędna $z$ ma postać $z+z'+xy'$.
Iloraz lewostronny $N_{\Z}\backslash N$ jest zwartą wiązką
okręgową nad $T^2$
o niezerowej liczbie Eulera. Jego grupa podstawowa
$N_{\Z}$ jest nilpotentna, lecz nieabelowa.

\textbf{Geometria $\mathrm{Sol}$.}
\index{geometrie Thurstona!Sol}
Niech $\R^2\rtimes\R$ ma mnożenie
\[
 (x,y,t)(x',y',t')=
 (x+e^t x',\;y+e^{-t}y',\;t+t')
\]
i lewostronnie niezmienniczą metrykę
\[
 g=e^{-2t}(\dd x)^2+e^{2t}(\dd y)^2+(\dd t)^2.
\]
Jedna pozioma oś rozszerza się przy wzroście $t$,
druga kurczy. Dla ortonormalnych kierunków $X,Y,T$
odpowiednio wzdłuż $x,y,t$ mamy
$K(X,Y)=1$ i $K(X,T)=K(Y,T)=-1$.
Składowa jedności grupy izometrii ma wymiar $3$.
Przykład zwartego ilorazu otrzymujemy z macierzy
\[
 A=\begin{pmatrix}2&1\\1&1\end{pmatrix}\in\mathrm{SL}(2,\Z).
\]
Jej wartości własne są dodatnie, wzajemnie odwrotne i różne od $1$.
Torusowa wiązka nad $S^1$ z monodromią $A$ ma grupę podstawową
$\Z^2\rtimes_A\Z$ i geometrię $\mathrm{Sol}$: w bazie
wektorów własnych $A$ jej działanie na $\R^2$ odpowiada
przesunięciu współrzędnej $t$ o $\log\lambda$, gdzie
$\lambda>1$ jest większą wartością własną.
Grupa ta jest rozwiązalna, ale nie jest wirtualnie nilpotentna.

\begin{uwaga}[Porównanie ośmiu modeli]
\label{uw:thurston-eight-comparison}
Trzy przestrzenie $S^3,\mathbb E^3,\mathbb H^3$ mają stałą
krzywiznę i sześciowymiarową grupę izometrii. Dwa iloczyny
mają płaszczyzny o krzywiźnie $0$ oraz odpowiednio $+1$ i $-1$.
Modele $\widetilde{\mathrm{SL}}(2,\R)$, $\mathrm{Nil}$
i $\mathrm{Sol}$ mają oba znaki krzywizny. Ich typowe ilorazy
odróżnia topologia: niezerowe skręcenie wiązki okręgowej
nad powierzchnią hiperboliczną, niezerowe skręcenie nad torusem
oraz hiperboliczna monodromia wiązki torusowej.
Sama informacja, że $K$ zmienia znak, nie wystarcza więc
do odróżnienia trzech ostatnich geometrii.
\end{uwaga}

\section{Włóknienia Seiferta i geometria}
\label{sec:thurston-seifert}

Rozmaitość Seiferta jest rozłożona na okręgi tak, że nad typowym
punktem powierzchni ilorazowej wygląda jak wiązka okręgowa.
Dokładniej, w przypadku orientowalnych włókien każde włókno ma
otoczenie będące ilorazem $D^2\times S^1$ przez działanie
cykliczne rzędu $m$ zadane wzorem
\[
 (z,u)\longmapsto
 (e^{2\pi i/m}z,e^{2\pi iq/m}u),\qquad \gcd(q,m)=1.
\]
Dla $m=1$ jest to zwykły iloczyn; dla $m>1$ włókno nad środkiem
dysku jest wyjątkowe. Iloraz bazy jest wtedy \emph{orbifoldem}
powierzchniowym z punktami stożkowymi rzędu $m$.
Dla orientowalnej bazy z punktami stożkowymi
rzędów $m_1,\ldots,m_r$ jej orbifoldowa charakterystyka Eulera to
\[
 \chi^{\mathrm{orb}}(B)=
 \chi(|B|)-\sum_{j=1}^r\left(1-\frac1{m_j}\right).
\]
Liczba Eulera włóknienia $e\in\Q$ mierzy jego skręcenie;
gdy nie ma włókien wyjątkowych, jest zwykłą całkowitą
liczbą Eulera wiązki okręgowej. Wystarcza nam rozróżnienie
$e=0$ i $e\ne0$, niezależne od konwencji znaku.
Na przykład wiązka iloczynowa $\Sigma_g\times S^1$ ma $e=0$,
natomiast $U\Sigma_g$ ma $e=\chi(\Sigma_g)\ne0$.

\begin{twierdzenie}[Geometria zamkniętej rozmaitości Seiferta]
\label{thm:thurston-seifert-table}
Załóżmy, że $M$ jest zamkniętą orientowalną rozmaitością
Seiferta z orientowalnym orbifoldem bazy $B$ oraz orientowalnymi
włóknami. Dla ustalonego włóknienia typ geometrii $M$
wyznaczają znak $\chi^{\mathrm{orb}}(B)$ i warunek $e=0$:
\[
\begin{array}{c|cc}
 & e=0 & e\ne0\\ \hline
\chi^{\mathrm{orb}}(B)>0 & S^2\times\R & S^3\\
\chi^{\mathrm{orb}}(B)=0 & \mathbb E^3 & \mathrm{Nil}\\
\chi^{\mathrm{orb}}(B)<0 & \mathbb H^2\times\R
 & \widetilde{\mathrm{SL}}(2,\R)
\end{array}
\]
Jest to klasyfikacja dla podanych założeń, przywołana z
\href{https://doi.org/10.1112/blms/15.5.401}
{P.\ Scotta, \emph{The Geometries of 3-Manifolds}, \S\,4,
tablica 4.1}. Przypadki z nieorientowalną bazą wymagają
odpowiednich wariantów i nie są objęte tą tabelą.
\end{twierdzenie}

Tabela wyjaśnia pary modeli o wspólnej bazie: nad powierzchnią
hiperboliczną brak skręcenia daje
$\mathbb H^2\times\R$, a skręcenie daje
$\widetilde{\mathrm{SL}}(2,\R)$. Nad torusem
otrzymujemy odpowiednio geometrię euklidesową i $\mathrm{Nil}$.
Geometrie hiperboliczna $\mathbb H^3$ i $\mathrm{Sol}$
nie występują w tej tabeli: ich zamknięte ilorazy nie są
rozmaitościami Seiferta. W szczególności wiązka torusowa
o hiperbolicznej monodromii z poprzedniej sekcji ma geometrię
$\mathrm{Sol}$, a nie $\mathrm{Nil}$.

\section{Twierdzenie geometryzacyjne}
\label{sec:thurston-theorem}

Sfera potrzebna do \emph{sumy spójnej} i torus potrzebny do
\emph{cięcia JSJ} pełnią różne role. Zamknięta trójrozmaitość
jest \emph{pierwsza}, jeśli z zapisu $M=M_1\#M_2$ wynika,
że jeden składnik jest $S^3$. Jest \emph{nierozkładalna},
jeśli każda gładko osadzona sfera $S^2$ ogranicza kulę $B^3$.
Zamknięte orientowalne składniki pierwsze są nierozkładalne
z wyjątkiem $S^2\times S^1$. Torus w nierozkładalnej
rozmaitości nazywamy \emph{nieściśliwym}, gdy inkluzja
$T^2\hookrightarrow M$ jest injektywna na $\pi_1$.
Równoważnie nie ma dysku ściskającego ten torus.
W tym sformułowaniu JSJ oznacza kanoniczne, po przyjęciu
zwykłych konwencji minimalności, cięcie wzdłuż rozłącznych
nieściśliwych torusów.

\begin{twierdzenie}[Geometryzacja Thurstona--Perelmana;
wynik zewnętrzny]
\label{thm:thurston-geometrization}
Niech $M$ będzie zamkniętą, spójną, orientowalną
trójrozmaitością. Istnieje rozkład na sumę spójną skończenie
wielu rozmaitości pierwszych, jednoznaczny z pominięciem
składników $S^3$ i kolejności. Każdy składnik
nierozkładalny można rozciąć wzdłuż skończonej rodziny
rozłącznych nieściśliwych torusów JSJ tak, że każda
powstała część jest jednego z dwóch rodzajów:
\begin{enumerate}[label=\textup{(\roman*)}]
\item jej wnętrze ma zupełną metrykę hiperboliczną
  o stałej krzywiźnie $-1$ i skończonej objętości;
\item jest rozmaitością Seiferta.
\end{enumerate}
Składnik pierwszy $S^2\times S^1$ ma geometrię
$S^2\times\R$ i nie wymaga cięcia JSJ.
\end{twierdzenie}

Jest to wygodna wersja geometryzacji po połączeniu
twierdzenia o rozkładzie pierwszym, rozkładu JSJ
i dowodu hipotezy Thurstona. Pełne dowody i dokładniejsze
warianty, także z butelkami Kleina, podają
\href{https://www.claymath.org/resource/the-geometrization-conjecture/}
{J.\ Morgan i G.\ Tian, \emph{The Geometrization Conjecture}}
oraz \href{https://arxiv.org/abs/0809.4040}
{ci sami autorzy, \emph{Completion of the Proof of the
Geometrization Conjecture}}. Trzeba odróżnić zdanie (i)
od twierdzenia, że każda część z brzegiem jest zwartym
ilorazem jednego z ośmiu modeli: jej \emph{wnętrze}
może być niezwartym ilorazem o skończonej objętości.
W przypadku (ii) część jest opisana przez włóknienie
Seiferta; tabelę z poprzedniej sekcji stosuje się bezpośrednio
do zamkniętych rozmaitości spełniających jej założenia.

\begin{przyklad}[Dwa proste wyniki cięcia]
\label{prz:thurston-cut}
Torus $T^3$ już sam jest zamkniętą rozmaitością geometryczną
typu euklidesowego;
nie trzeba go przecinać, mimo że zawiera wiele nieściśliwych
torusów. Dla rozmaitości sklejonej z dwóch części Seiferta
wzdłuż torusa włókna okręgowe mogą spotykać się pod różnymi
nachyleniami. Wtedy całość nie musi mieć jednej geometrii
modelowej, choć obie części są typu Seiferta. Rozkład JSJ
zachowuje informację o miejscu sklejenia.
\end{przyklad}

\section{Droga od potoku Ricciego do geometryzacji}
\label{sec:thurston-ricci-outline}

Rozdział~\ref{ch:potok-ricciego} wprowadza równanie
$\partial_tg=-2\operatorname{Ric}(g)$ i dowodzi istnienia
rozwiązania przez krótki czas. Aby uzyskać geometryzację,
potrzebne są cztery dalsze, głębokie kroki.

Po pierwsze, gdy krzywizna rośnie i gładki potok zbliża się
do osobliwości, analizuje się powiększone otoczenia punktów
o dużej krzywiźnie. Niezapadanie objętości i opis otoczeń
kanonicznych pozwalają rozpoznać cienkie szyjki oraz ich
zamknięte końce. Po drugie, chirurgia przecina wybrane szyjki,
zakleja końce i pozwala prowadzić potok dalej, jednocześnie
kontrolując zmianę topologii. Konstrukcję i jej analityczne
założenia opisuje
\href{https://arxiv.org/abs/math/0211159}
{G.\ Perelman, \emph{The entropy formula for the Ricci flow and
its geometric applications}} oraz
\href{https://arxiv.org/abs/math/0303109}
{G.\ Perelman, \emph{Ricci flow with surgery on three-manifolds}}.

Po trzecie, dla długiego czasu odpowiednio przeskalowany
potok rozdziela obszary \emph{grube}, z których powstają
hiperboliczne części o skończonej objętości, od obszarów
\emph{cienkich}, lokalnie zapadających się objętościowo.
Po czwarte, analiza tych ostatnich prowadzi do rozmaitości
grafowych, zbudowanych z części Seiferta. Ten ostatni krok
jest istotną częścią dowodu Morgana i Tiana, a nie
bezpośrednią konsekwencją samego równania potoku.

W przypadku jednospójnym istnieje krótsza droga do wniosku
topologicznego. Twierdzenie Perelmana o wygaśnięciu mówi:
jeśli zamknięta orientowalna trójrozmaitość ma rozkład
pierwszy bez składników asferycznych, to skonstruowany
potok z chirurgią wygasa w skończonym czasie dla każdej
gładkiej metryki początkowej
(\href{https://arxiv.org/abs/math/0307245}
{G.\ Perelman, \emph{Finite extinction time for the solutions
to the Ricci flow on certain three-manifolds}}).
Dla rozmaitości jednospójnej warunek ten jest spełniony:
grupa podstawowa sumy spójnej jest iloczynem wolnym grup
składników, a zamknięta asferyczna trójrozmaitość nie jest
jednospójna. Gdyby była, byłaby ściągalna, czemu przeczy
jej niezerowa najwyższa homologia ze współczynnikami $\Z/2$.
Opis topologii usuwanych składników pokazuje
wtedy, że rozmaitość jest sumą spójną sferycznych form
przestrzennych i składników $S^2\times S^1$.
Jej grupa podstawowa jest iloczynem wolnym odpowiednich
grup skończonych i kopii $\Z$. Jeśli jest trywialna,
wszystkie te czynniki są trywialne; pozostaje $S^3$.
Ścisłe przejście przez chirurgię i klasyfikację usuwanych
składników znajduje się w
\href{https://www.claymath.org/resource/ricci-flow-and-the-poincare-conjecture/}
{J.\ Morgana i G.\ Tiana, \emph{Ricci Flow and the Poincaré
Conjecture}}. Sam wniosek i jego miejsce wśród wersji
twierdzenia Poincarégo omawia Rozdział~\ref{ch:poincare-trzy}.

Powyższy łańcuch jest \emph{opisem idei dowodu}, a nie
dowodem twierdzenia geometryzacyjnego ani twierdzenia
Poincarégo. Analiza otoczeń osobliwych, chirurgii i zapadania
się wymaga osobnych twierdzeń, których założenia i dowody
znajdują się w podanych źródłach.
